\section{Preparation of the two-endpoint statuses}
\label{sec:preparation}

The scalar recipes involve the endpoint functional
\(p_i=u_{i'}U_i\), where \(i'\) is the other endpoint of the unit.
We distinguish laws on which the values \(p_i(x)\) have a common description in
terms of the effective-space evaluations, the role bit, and the
global key. Such a description will be called a prediction.
For a predicted law, we combine the phase estimates with two
properties of predictions to obtain the compatible tables stated
below. If no prediction exists, the first peeling already supplies
the preparation we need. The scalar recipes are constructed in
\Cref{sec:status-span} for the positive-mass alternatives, and
directly in \Cref{sec:rare-overlap} for the inverse-polynomial
alternative.

We work along a sequence of admissible unit laws that would violate
\Cref{thm:raw-supersaturation}, and may pass to subsequences.
Perform the first peeling and the deletion in
\Cref{lem:raw-intersections}. The effective spaces at this stage
are denoted by \(C_i^{\mathrm{old}}\). For a tag \(l\) and a
diagonal bit \(q\), let \(\Omega_{l,q,n}\) be the uniform law on
the global plus and minus vector slots of the star of \(l\),
with paired diagonal \(q\) in every component. Requiring its
individual vectors to be nonzero changes it by exponentially
small total variation.

\subsection{Common predictions and the preparation statement}

\begin{definition}[Common prediction]
\label{def:prediction}
A prediction along a subsequence consists of common binary
functions
\[
 f_{i,n}(l,q,\cdot):\Omega_{l,q,n}\longrightarrow\F,
 \qquad i=1,2,
\]
a set of units of mass at least \(10^{-6}\), and, on every
unit in this set, choices
\[
 \lambda_i\in C_i^{\mathrm{old}},\qquad
 \alpha_i\in\F,\qquad
 {\cal E}_i\subseteq\F^b,\quad |{\cal E}_i|\le B_*,
\]
such that the following uniform statement holds. There is
\(\varepsilon_n\to0\) for which, at both endpoints, at every
label outside \({\cal E}_i\), at every tag, and in every
permitted flavor, a uniform parameter atom \(x\) satisfies
\begin{equation}
 p_i(x)=T(\lambda_i,x)+\alpha_i a(x)
          +f_{i,n}(l,q,\operatorname{key}(x))
 \label{eq:prediction}
\end{equation}
except with probability at most \(\varepsilon_n\).
The functions \(f_{i,n}\) are the same for all retained units.
The marks \(\lambda_i,\alpha_i,{\cal E}_i\) may depend on
the whole unit.
\end{definition}

The errors refer to uniform parameters in each flavor, before
conditioning on a tester bit. Conditioning on a possible fixed
tester value only multiplies the error by a fixed constant.
There are finitely many labels, tags and flavors, with their
numbers fixed independently of \(n\).

If a prediction exists, retain its marked population and
pigeonhole the values of the two \(\alpha_i\)'s, of
\[
 c=a(\lambda_1)+a(\lambda_2),
\]
and of the indicator \(\Delta=0\), where
\(\Delta=U_1\lambda_1+U_2\lambda_2\).
This loses a factor at most \(16\), leaving mass at least
\(6.25\cdot10^{-8}>10^{-12}\).
Pass to a subsequence on which every reference error probability
of every binary combination of \(f_1,f_2,q\) has a limit,
for each tag and bit. Only finitely many comparisons are needed.

Identify two sequences of common functions when their disagreement
probability tends to zero on every \(\Omega_{l,q,n}\), and then
quotient by the span of the common diagonal-bit function \(q\).
Thus two predictors have the same class when their difference
agrees asymptotically with \(\gamma q\), for a single
\(\gamma\in\F\) common to all tag and bit spaces.
Write \([f]\) for this class and put
\begin{equation}
 d_i=1+\alpha_i,\qquad e_i=(d_i,[f_i]).
 \label{eq:status-class}
\end{equation}
The \(e_i\) are \emph{prediction classes}, with support
\(\{i:e_i\ne0\}\). They are common to all retained units,
although the marked \(\lambda_i\)'s can vary. A query's
\emph{numerical unary status}, by contrast, records the scalar
outcomes specified in \Cref{sec:collision}.

For an admissible small table, contractions on an old
\(\lambda_i\) are computed from the table. They remain defined
after pins are enlarged, since the old effective spaces are
retained inside the new ones. We require the following phase
compatibility conditions:
\begin{equation}
\begin{array}{ll}
\displaystyle
u_{B,S^\circ}(\Delta_A)+u_{A,S^\circ}(\Delta_B)
       =|S^\circ|d,
&
\begin{gathered}
\text{if }S^\circ=\{i:e_i\ne0\}\ne\varnothing,\\[-1mm]
\text{and all nonzero }e_i\text{ coincide},\\[-1mm]
d=\text{their first bit};
\end{gathered}
\\[4mm]
u_{B,j}(\Delta_A)=c,\quad u_{A,i}(\Delta_B)=c
\quad(i,j=1,2),
&\text{if }e_1=e_2=0.
\end{array}
\label{eq:phase-compatibility}
\end{equation}
There is no additional phase condition in the other cases.
When testing a table, the \(u\)'s in this display mean its
contractions, rather than contractions of actual orientations.
These equations remove the possible dual obstructions to the
scalar recipes, as proved in \Cref{lem:status-predicted}.

\begin{proposition}[Status preparation]
\label{prop:status-preparation}
Along a subsequence, the unit law can be prepared with final
pin dimension at most \(K\) and the final per-leaf bounds of
\Cref{lem:peeling}, at relative mass cost \(2^{-o(N)}\), so
that one of the following descriptions applies.
\begin{enumerate}[label=\textup{(\roman*)}]
\item There is no prediction as in \Cref{def:prediction}
on the first-peeled law. The effective spaces are still the
first ones. For independent units, the mass admitting an
admissible injecting small table is greater than \(.99\).
\item A prediction with fixed \(\alpha_i,c\) and common
classes \(e_i\) is retained. The old marks and exceptions
are preserved, and \Cref{eq:prediction} has uniform \(o(1)\)
error on every retained unit. The mass admitting an
admissible injecting table satisfying
\Cref{eq:phase-compatibility} is bounded below by a positive
constant.
\item Both prediction classes are zero, and every retained
unit satisfies the exact identities
\[
 p_i=T(\lambda_i,\cdot)+a
 \quad\text{on all of }\cX,\qquad i=1,2.
\]
There is a set of leaf pairs of mass at least
\(c_2N^{-4}\) under the final leaf mixture, for a fixed
\(c_2>0\), with phase compatibility determined by the
two leaves. For each such pair one can choose a specific
injecting compatible table whose two separate orientation
admissibility filters have masses at least \(a_0>0\).
The constants \(c_2,a_0\) are independent of \(n\).
\end{enumerate}
In every case the mixed joint density is at most
\(2^{(D+o(1))N}\mu^2\), and its individual marginals are at
most \(M2^{o(N)}\mu\).
Relative to the stored leaf bases and pin images, all numerical
table formats, unary statuses, old-mark coordinates and exception-set
descriptions have bounded length in \(n\). The extra exact pins
are used only in case \textup{(iii)}.
\end{proposition}

We first prove two facts about predictions. Vanishing prediction
classes give exact identities on \(\cX\), while \(\Delta=0\)
forces \(\alpha_1=\alpha_2\). The latter fact makes the required
phase equation attainable when the tensor difference vanishes.
After these proofs, the phase estimates of \Cref{sec:phases}
give the preparation alternatives.

\subsection{Exact identities for vanishing prediction classes}

\begin{lemma}[Exactification of vanishing prediction classes]
\label{lem:empty-status}
If \(e_1=e_2=0\), deleting \(o(1)\) unit mass makes
\begin{equation}
 p_i=T(\lambda_i,\cdot)+a
 \quad\text{on all of }\cX,\qquad i=1,2,
 \label{eq:exact-prediction}
\end{equation}
an exact identity on every retained unit.
\end{lemma}

\begin{proof}
Vanishing prediction classes give \(\alpha_i=1\) and
\(f_i\simeq\gamma_iq\) for a fixed \(\gamma_i\in\F\).
At any fixed tag, label and parameter point, the raw key law
is the single-key reference up to an exponentially small
injectivity error, by \Cref{lem:frame-images}.
The individual marginal density bound, which is a fixed
constant at this stage, transfers the reference disagreement
of \(f_i\) and \(\gamma_iq\) to an average \(o(1)\) parameter
error on the marked population.
There are only finitely many endpoint, tag and label choices.
Markov's inequality therefore allows deletion of \(o(1)\)
unit mass so that this error tends to zero uniformly on
the remaining units and choices.
Combining it with \Cref{eq:prediction} gives, at every
nonexceptional label,
\[
 p_i(x)+T(\lambda_i,x)+a(x)+\gamma_iq(x)=0
\]
with parameter error tending to zero in the generic flavor.

For fixed label and tag the left side is a Boolean polynomial
of base degree at most two.  A nonzero Boolean polynomial of
degree at most \(k\) has nonzero probability at least \(2^{-k}\)
under uniform bits.  For completeness, this follows by induction:
choose a coordinate with a nonzero derivative; the derivative
has degree at most \(k-1\), and whenever it is nonzero at least
one of the two values on that coordinate is nonzero.
The constant nonzero polynomial is the initial case.
Thus the displayed relation, once its error is below \(1/4\),
is exact for every allowed base value.

For any fixed allowed base value its dependence on the selector
has degree at most \(2j_*\).  It vanishes at every selector
outside a set of size at most \(B_*\).
The choice of \(b\) gives
\(B_*<2^{b-2j_*}\), so the same polynomial bound forces it
to vanish at the excluded labels also.

Use one shared-only base point with \(\eta_S=1\), the same at
every tag.  The sum of its \(g\) star atoms is zero in the cut
space.  Summing the exact relation over these atoms gives
\(0=\gamma_i g=\gamma_i\), because \(g\) is odd and
the first three terms are linear on \(\cX\).
Consequently \(\gamma_i=0\).
The point atoms span the cut space by its definition, proving
\Cref{eq:exact-prediction}.
\end{proof}

\subsection{The zero-difference case}

We next show that \(\Delta=0\) forces \(\alpha_1=\alpha_2\).
The proof compares the common predictors on affine parameter slices.
The predictors need not be polynomial, so we first establish the
comparison for arbitrary common functions of the displayed image
columns. It does not apply to arbitrary orientation-dependent tests.

\begin{lemma}[Affine-slice comparison]
\label{lem:synthetic-slice}
Fix a selector, a flavor, and a fixed number \(d\) of affine
parameter variables.  Write a random affine coefficient map as
\[
 C(t,u)=p_s\otimes(t,z_0t+Lu),
 \qquad (t,u)\in\F\oplus\F^d.
 \]
Assume the flavor retains an \(n\)-bit \(Z\)-block absent from
\(E\).  Suppose we condition the affine coefficients on
\(C\trans EC=G\), where \(G\) is fixed, this event has probability
at least a fixed \(\beta>0\), and it depends on no \(Z\)
coefficient.
Let \({\cal T}_n\) be any common test of absolute value at most
one on the image column arrays \((P_eC,Q_eC)_e\), using a fixed
set of components.  Its empirical average over the conditioned
affine coefficients, at a raw orientation, has variance
\[
 O(2^{-N/2}+2^{2d+1-n})
 \]
around its mean under uniform injective image column arrays
with internal Gram \(G\).
The constants may depend on the fixed column counts and
\(\beta\), but the estimate is uniform over the tests
\({\cal T}_n\).
\end{lemma}

\begin{proof}
The \(Z\)-coefficients remain independent uniform bits after
the conditioning.  The columns of \(C\) are independent
except with probability \(O(2^{d-n})\).
For two independent conditioned maps \(C,C'\), a relation
\[
 C(a,u)+C'(a',u')=0
 \]
has \(a=a'\), by its constant coordinate.  Its \(Z\)-part is
\[
 a(z_0+z_0')+L_Zu+L'_Zu'=0.
 \]
The matrix \([z_0+z_0',L_Z,L'_Z]\) is a uniform
\(n\times(2d+1)\) bit matrix.  A union over its nonzero
coefficient vectors bounds the failure of combined independence
by \((2^{2d+1}-1)2^{-n}\).
In particular, the two translation columns do not share a
forced nominal constant direction.

For fixed independent nominal columns, the restricted-image
transitivity in \Cref{lem:frame-images} gives the uniform
injective law with its specified Gram.  One batch consequently
has the claimed reference law, independently of the remaining
nominal coefficients.
For two batches, compared with two independent copies of this
reference law, the actual images have finitely many additional
mutual Gram entries prescribed.  Joint injectivity contributes
only \(O(2^{-N+O(d)})\) to the comparison: under independent
uniform vector slots, a dependence among a fixed number of
columns has this probability, and the fixed Gram conditioning
has bounded reciprocal probability.

Expand the mutual Gram indicator in binary characters, collapsing
any repeated identical entries.  A nonzero character has a
nonzero coefficient matrix on at least one interaction between
the two batches.  More explicitly the coefficient blocks pairing
first-batch plus vectors with second-batch minus vectors and
first-batch minus vectors with second-batch plus vectors give
cross-batch bit rank
\(N(\rank A+\rank B)\ge N\).
The internal Gram constraints, individual injectivity indicators,
and the two arbitrary tests are separate bounded weights on
the two batches; their normalization factors are fixed.
\Cref{lem:walsh} bounds every nontrivial character contribution
by \(O(2^{-N/2})\).  The same computation without the tests
gives the normalization of the mutual Gram conditioning.
There are only a fixed number of characters.  Thus the
conditional two-batch product expectation differs from the
product of the one-batch expectations by \(O(2^{-N/2})\),
uniformly over the nominal mutual Gram values.

Averaging the coefficients, and including the probability of
nominal dependence, proves the asserted second-moment estimate.
The one-batch mean differs from the reference mean only by
the already bounded nominal-dependence error.
\end{proof}

\begin{lemma}[Equal coefficients when the tensor difference vanishes]
\label{lem:zero-delta}
On a predicted branch with \(\Delta=0\) throughout,
\(\alpha_1=\alpha_2\).
\end{lemma}

\begin{proof}
Suppose otherwise.  By \Cref{lem:raw-intersections}, the two
canonical sparse representations \(w^{(i)}=(w_l^{(i)})_l\)
of the marked \(\lambda_i\)'s have the same value of \(\chi_*\).
Each has at most \(4D/g\) nonzero tag entries.
As \(\alpha_1\ne\alpha_2\), on every unit exactly one endpoint
has \(\chi_*(w^{(i)})+\alpha_i=1\).
Fix an endpoint \(i\) on a positive fraction of the units.
Pigeonhole the set
\[
 T_0=\{t:\eta_S(w_t^{(i)})=1\},
 \qquad |T_0|\le4D/g,
 \]
and choose a single selector outside the stored exceptions
on a positive fraction of this remaining population.
The latter is possible by averaging over selectors, since
\(B_*<2^{b-1}\).
All these restrictions have fixed positive mass in \(n\);
their endpoint marginal is at most \(L\mu\) for a fixed \(L\).
No restriction made solely for this contradiction will be
used later to choose a channel-size threshold.

At each tag \(l\), use the pure flavor retaining exactly the
allowed \(O_{d,t}\) blocks with \(t\notin T_0\), setting all
\(S\) bits to zero and retaining the \(\#,Z\) blocks.
For a pure tag atom, \(a(x)=q(x)\) and every \(b_t(x)=0\).
The coefficient of its tester \(\eta_{d,t}\) in
\(T^0(\lambda_i,x)+\alpha_i a(x)\) is
\[
 a(\lambda_i)+b_t(\lambda_i)+\alpha_i
 =\chi_*(w^{(i)})+\eta_S(w_t^{(i)})+\alpha_i=1
 \]
on the retained \(t\)'s.  Hence \Cref{eq:prediction} gives
\begin{equation}
 f_i(l,q,\operatorname{key}(x))+q
 =p_i(x)+T^1(\lambda_i,x)
 \label{eq:pure-prediction}
\end{equation}
with uniform \(o(1)\) parameter error.

The right side of \Cref{eq:pure-prediction} is a sum of at
most
\begin{equation}
 s_1=(g-1)(h+2JK_1)
 \label{eq:synthetic-product-count}
\end{equation}
products of affine functions of the base variables.
Indeed \(p_i\) contributes at most \(h\) products in each
of the \(g-1\) star components.  Factoring the old
\(\lambda_i\)'s, whose total component rank is at most
\(K_1\), gives at most \(2J(g-1)K_1\) products for \(T^1\).
The selected pure tester retains at least
\begin{equation}
 \frac{g-1}{2}(g-4D/g)r_0>2s_1+1
 \label{eq:synthetic-tester-room}
\end{equation}
disjoint ordered bit pairs.  To check the last inequality,
use \(D=4000g\) and \(h=1000r_0\).  After division by
\((g-1)r_0\), the left side is \((g-16000)/2\) and the
right side is at most \(2000+4JK_1/r_0+o(1)\).
The fixed value \(g=10^9+1\) and the later choice of
sufficiently large \(r_0\) make the inequality hold.

We will compare \(g\) matrices of cross coefficients, each of rank
at most \(2s_1\), whose sum is an identity matrix of size
\(2gs_1+1\). To construct them, place the common predictor on one
array of synthetic keys. Shared-only parity will determine the sum,
while its restrictions to small pure-flavor slices will give the
individual rank bounds.

Put \(m_1=2gs_1+1\), and introduce synthetic variables
\(\xi,\upsilon\in\F^{m_1}\) and
\[
 q_*=\sum_{j=1}^{m_1}\xi_j\upsilon_j.
 \]
For every component independently, choose uniform injective
maps of the affine coefficient columns
\((1,\xi,\upsilon)\) into the plus and minus ambient spaces
with full paired Gram
\[
 \pair{\text{plus}(1,\xi,\upsilon)}
       {\text{minus}(1,\xi',\upsilon')}
 =\sum_j\xi_j\upsilon'_j.
 \]
Such maps exist for large \(n\), since the number of columns
is fixed.  Their point evaluations form one synthetic array
of global component keys; all tag stars use the corresponding
entries of this same array.

We next prove two properties of this array with probability
tending to one.

\medskip\noindent
\emph{Shared-only parity.}
In a local shared-only test use the fixed selector and the
same point at every tag.  The sum of the tag atoms is zero
in \(\cX\).  Summing \Cref{eq:prediction} therefore gives
\[
 \sum_l f_i(l,q,\operatorname{key}_l)=0
 \]
except on \(o(1)\) parameter mass, uniformly on the retained
units.  Both tester bits have probability at least \(1/4\):
the bias of the sum of \(r_0\) independent bit-pair products
is \(2^{-r_0}\).

Apply \Cref{lem:synthetic-slice} with no affine variables
(\(d=0\)), the shared-only flavor, and internal Gram \((q)\).
The Gram event depends only on the fixed tester coordinates
and has probability at least \(1/4\).  The common test is
the displayed parity predicate of the global keys in all
components.  Its empirical frequency under \(\mu\) concentrates
at its global reference probability.  The endpoint marginal
bound \(L\mu\), and its frequency \(1-o(1)\) on the retained
population, force that reference probability to tend to one.
At every point of the synthetic array the point-key marginal
is this reference, up to the exponentially small nonzero-vector
conditioning.  There are a fixed finite number of points.
Thus, with probability tending to one,
\begin{equation}
 \sum_l f_i(l,q_*,\operatorname{key}_l)=0
 \quad\text{at every synthetic point}.
 \label{eq:synthetic-parity}
\end{equation}

\medskip\noindent
\emph{Affine-product representations on every small slice.}
Fix sets \(I,J\subseteq\{1,\ldots,m_1\}\), each of size at
most \(2s_1+1\), and set all \(\xi\) and \(\upsilon\)
coordinates outside \(I,J\), respectively, to zero.
There are \(d=|I|+|J|\le4s_1+2\) remaining variables.
At a given tag, map this slice into its selected local pure
flavor by independent uniform affine functions for every
retained base bit, including independent translations.
Condition their full oriented Gram to equal the synthetic
slice Gram.

This conditioning has a fixed positive probability.
At the fixed selector \(E^\#\) vanishes on pairs.
There are at most \(2s_1+1\) active terms
\(\xi_j\upsilon'_j\) in the target Gram, so
\Cref{eq:synthetic-tester-room} lets us route them into
distinct retained ordered tester pairs, and set all other
tester affine coefficients to zero.  If \(t_0\) scalar
tester coordinates are prescribed, this assignment alone
has probability \(2^{-t_0(d+1)}>0\), independent of \(n\).
The \(Z\)-affine coefficients remain iid.
At every fixed slice point the unconditioned base input is
uniform; after conditioning, its density is bounded by the
fixed reciprocal Gram probability.

Consequently \Cref{eq:pure-prediction} holds simultaneously
at all points of the slice except on \(o(1)\) of these
parametrizations, uniformly on the marked population.
Its right side pulls back to a sum of at most \(s_1\)
affine products.  Consider only the event
\begin{quote}
the array of values of the common function \(f_i+q_*\) on
this specified slice admits such an affine-product
representation.
\end{quote}
This is a common test of the image column arrays.  It contains
neither the mark \(\lambda_i\) nor any coefficients of the
unit-specific representation.  Its empirical success
probability is \(1-o(1)\) on the retained population.
\Cref{lem:synthetic-slice} and the marginal bound \(L\mu\)
therefore force its probability under the synthetic slice
law to tend to one.  That law is also the marginal of the
corresponding slice in the full synthetic array, by restricted
frame transitivity.

There are only finitely many choices of tag and slice, with
the number fixed independently of \(n\).  We may therefore
require all these events and \Cref{eq:synthetic-parity}
simultaneously.  Notice that we embedded only the small
slices locally; no local realization of the entire synthetic
tester was required.

Choose one synthetic array on which all these properties hold,
and define
\[
 F_l(\xi,\upsilon)
   =f_i(l,q_*,\operatorname{key}_l)+q_*.
 \]
For each tag define its Boolean cross-coefficient matrix by
\[
 C_l[j,k]=F_l(e_j,e_k)+F_l(e_j,0)
                 +F_l(0,e_k)+F_l(0,0).
 \]
On a coordinate slice where \(F_l\) is a sum of \(s_1\)
affine products, this is the matrix of \(\xi_j\upsilon_k\)
coefficients.  One affine product contributes a sum of two
rank-one matrices to this cross matrix.  Every
\((2s_1+1)\)-square minor of \(C_l\) is therefore singular,
so \(\rank C_l\le2s_1\).
On the other hand \Cref{eq:synthetic-parity}, together with
odd \(g\), gives \(\sum_lF_l=q_*\) and hence
\[
 \sum_l C_l=I_{m_1}.
 \]
Rank subadditivity would imply
\(m_1\le\sum_l\rank C_l\le2gs_1\), contradicting
\(m_1=2gs_1+1\).  This proves \(\alpha_1=\alpha_2\).
\end{proof}

\subsection{Proof of the preparation statement}

\begin{proof}[Proof of \Cref{prop:status-preparation}]
If no prediction exists along any subsequence, retain the
first-peeling law.  Apply \Cref{lem:injection} to the
unconditional pair event, using the first pins and the final
budget \(K\) as an upper bound. The failure tolerance here and
in the remaining cases is a fixed numerical constant chosen before
\(h\). With a sufficiently small failure tolerance it gives
injection probability greater
than \(.99\).  The actual small table is admissible, so this
proves (i).  Negligible later deletions cannot create a
prediction on a mass bounded above \(10^{-6}\) by a fixed
positive amount: transfer its uniformly valid marks and
functions to the original law and multiply its mass by
the retained proportion.  This is the form of the
no-prediction conclusion used subsequently.

Otherwise make the prediction and the finite initial
pigeonholing described after \Cref{def:prediction}.
If both prediction classes vanish, perform the \(o(1)\) deletion
in \Cref{lem:empty-status} now.  Thus all later refinements
start with the exact identity in that case.
For the moment test attainability of
\Cref{eq:phase-compatibility} on the actual table.

\paragraph{Vanishing tensor difference.}
Suppose first that \(\Delta=0\).
If there is exactly one active class and its first bit
were \(d=1\), its endpoint would have \(\alpha_i=0\),
whereas the inactive endpoint would have \(\alpha_{i'}=1\).
This contradicts \Cref{lem:zero-delta}.
With two equal active classes the right side of
\Cref{eq:phase-compatibility} is \(2d=0\).
Thus all nonempty-support conditions are automatic.
If both classes vanish and \(c=1\), the equality
\(U_1\lambda_1=U_2\lambda_2\), together with
\Cref{eq:exact-prediction}, gives the two affine-gradient
equations and the required unequal \(a\)-values for an
internal hole.  This is impossible for a unit.
Consequently \(c=0\), and the empty-support phase conditions
are automatic too.  Unconditional injection gives (ii).

\paragraph{Mixing of the four phase bits.}
Now suppose \(\Delta\ne0\), and apply
\Cref{thm:phase-alternative}.
If its mixing alternative holds, any specified pattern of
the four cross-phase bits has probability at least
\[
 \frac{1-15(.005)}{16}>.05
 \]
by binary Fourier inversion.  Choose a pattern satisfying
\Cref{eq:phase-compatibility}, if any condition is required.
Unconditional injection fails on an arbitrarily small
fixed pair mass by \Cref{lem:injection}.  Taking that mass
below \(.01\) leaves a positive mass of actual tables with
all required properties, proving (ii).

\paragraph{A tiny cover with nonempty support.}
It remains to treat the tiny-cover alternative.
If there is no extra phase condition, injection alone
suffices.  If there is one active class, the only
condition is one singleton phase equality, which has
probability \(1/2+2^{-\Omega(N)}\).  Again unconditional
injection leaves positive mass.
If both classes are nonzero and equal, the required bit
is
\[
 F(A,B)+F(B,A)=0,
 \qquad F(A,B)=u_{B,\{1,2\}}(\Delta_A).
 \]
Pass to a subsequence on which its probability has a
limit.  If that limit is positive, use the accepting
family in \Cref{rem:tiny-accepting-family} and the relative
form of \Cref{lem:injection}; a positive mass remains.
If the limit is zero, part (a) of \Cref{lem:tiny-cover}
retains a fiber of mass \(2^{-o(N)}\) on which the
condition is automatic.  Its individual marginals have
at most the allowed \(2^{o(N)}\) inflation.
Unconditional injection on that fiber gives a positive
mass of suitable tables in the normalized law.
All these cases give (ii).

\paragraph{A tiny cover with vanishing classes.}
Finally suppose both prediction classes vanish.
For any unit,
\[
 F(A,A)=u_1(U_2\lambda_2)+u_2(U_1\lambda_1)
       =p_2(\lambda_2)+p_1(\lambda_1)=0.
 \]
The same-endpoint contractions vanish by \(u_iU_i=0\),
and \Cref{eq:exact-prediction} gives
\[
 p_i(\lambda_i)=T(\lambda_i,\lambda_i)+a(\lambda_i)
              =a(\lambda_i)+a(\lambda_i)=0.
 \]
Parts (b) and (c) of \Cref{lem:tiny-cover} now give
four-phase compatibility on \(\Omega(N^{-4})\) leaf-pair
mass, determined by those leaves after at most \(u_0\)
additional pins and the second peeling.
The exponential deletions leave the same polynomial
order of mass.  Apply \Cref{lem:injection} to this
compatibility event, using the family in
\Cref{rem:tiny-accepting-family}.  Its relative failure
can be made less than \(.001\).

For clarity, disintegrate this last assertion over the
independent leaf pair.  A positive fixed fraction of the
compatible leaf-pair mass has conditional injection
probability at least \(1/2\); indeed a conditional
failure larger than \(1/2\) can account for at most twice
the total failure mass.  There are at most a fixed
number \(Q\) of numerical small tables in any leaf-pair
format, since all table-coordinate dimensions are bounded
by constants in \(n\).  For each of these good pairs,
some actual injecting table value occurs on conditional
pair mass at least \(1/(2Q)\).  This table is compatible:
the old tensor supports are pinned, so all its contractions
on the old marks agree with the fixed actual values.

For that table, admissibility is the intersection of one
filter on the first orientation and one on the second,
because pin images are fixed in the leaves.  Their
conditional product mass is at least \(1/(2Q)\), since
it includes the event where the actual table equals the
chosen one.  Each separate filter therefore has mass
at least \(a_0=1/(2Q)\).  Choosing such a table for each
good leaf pair proves (iii), without pigeonholing over
the number of leaves or ambient image values.

In the positive-mass cases the existence statement also
has the required finite-data form.  Table entries and
pin/projection relations are bounded numerical data.
The old marks can be recorded by their coordinates in
the old effective bases, and exception labels have
bounded descriptions.  For a fixed leaf pair the
admissibility checks are separate unary orientation
flags.  The phase contractions are computed from the
table and these bounded mark coordinates.  Thus the
actual-table success is detected among a finite list of
table/metadata alternatives; no pairwise conditioning
of the two unit draws is needed.

\paragraph{Retained masses and leaf bounds.}
We finish by recording the density and peeling effects.
All global restrictions above have mass at least
\(2^{-o(N)}\); all except the alternating fiber have
fixed positive mass, apart from deletions tending to
zero.  Therefore the mixed law has the density caps
in the statement.  If a retained portion of a first
leaf has relative mass less than \(2^{-\zeta N}\),
discard that portion.  Summed over the old leaf
weights, these losses have exponentially small mass,
also after division by the global retained mass
\(2^{-o(N)}\).  On the other leaves the joint
exact-image bound weakens from \(1-\zeta\) to at most
\(1-2\zeta\), as in \Cref{lem:peeling}.
Use its second-peeling conclusion only in case (iii),
where the new image rank is within \(u_0\).
This leaves total rank at most \(K\) and all stated
conditional leaf bounds.  The uniform unitwise errors
of \Cref{eq:prediction} survive every restriction;
when both classes vanish, the identity is already exact.
This proves all assertions.
\end{proof}

\begin{remark}[No feedback from later accuracies]
\label{rem:preparation-accuracies}
The injection tolerances used in this section are fixed numerical
constants: they are chosen before \(h\), together with the early
phase thresholds.  Injection is established before splitting
success among numerical tables.  Thus the potentially much smaller
flag lower bound \(a_0=1/(2Q)\), which may depend on \(h\), never
becomes an input to the choice of \(h\).
The errors in the affine-slice comparison and in subsequent
fixed-test limiting arguments are handled by increasing \(n\).
Their accuracies, slice counts, partition sizes, and flag-mass
thresholds do not require any alteration of the already fixed
channel size or of the early pin budgets.
\end{remark}
